\subsection{Trilhas de Leitura --- Projeto Hermathena}

O Projeto Hermathena pode ser explorado por diferentes trilhas de leitura, dependendo do perfil e interesse do leitor. Cada trilha seleciona os documentos mais relevantes para um tipo de público específico.


\bigskip
\centerline{\rule{13cm}{0.4pt}}
\bigskip

\subsubsection{Trilha 1 --- Visão Geral Institucional}

\textbf{Para:} Gestores, parceiros e financiadores.\newline
\textbf{Objetivo:} Compreender o escopo, a estrutura e os recursos do projeto.

\bigskip\noindent
\begin{tabular}{p{\dimexpr 0.333\linewidth-2\tabcolsep}p{\dimexpr 0.333\linewidth-2\tabcolsep}p{\dimexpr 0.333\linewidth-2\tabcolsep}}
\toprule
Ordem & Documento & O que encontrar \\
\hline
1 & \texttt{01\_motivacao.md} & Origem do projeto, missão, pilares estratégicos e ODS \\
2 & \texttt{02\_contexto-bingo.md} & O que é o BINGO, missão científica e linha do tempo \\
3 & \texttt{01\_estrutura.md} & Organização em núcleos, funções e dinâmica de trabalho \\
4 & \texttt{cronograma.md} & Cronograma de implantação por fases \\
5 & \texttt{orcamento.md} & Orçamento consolidado por núcleo e fontes de financiamento \\
6 & \texttt{parceiros.md} & Parceiros estratégicos, editais e financiadores \\
\bottomrule
\end{tabular}

\bigskip


\bigskip
\centerline{\rule{13cm}{0.4pt}}
\bigskip

\subsubsection{Trilha 2 --- Estratégia de Comunicação Científica}

\textbf{Para:} Comunicadores, jornalistas e educadores.\newline
\textbf{Objetivo:} Entender a abordagem metodológica e os canais de comunicação.

\bigskip\noindent
\begin{tabular}{p{\dimexpr 0.333\linewidth-2\tabcolsep}p{\dimexpr 0.333\linewidth-2\tabcolsep}p{\dimexpr 0.333\linewidth-2\tabcolsep}}
\toprule
Ordem & Documento & O que encontrar \\
\hline
1 & \texttt{01\_educomunicacao.md} & Fundamentação teórica: educomunicação, comunicação pública da ciência \\
2 & \texttt{ecossistema.md} & Visão integrada dos núcleos e seus fluxos \\
3 & \texttt{02\_Eixos\_Tematicos.md} & Eixos temáticos de conteúdo (Tengre, Shiva Nataraja, Sun Wukong, etc.) \\
4 & \texttt{04\_nucleo-via-lactea.md} & Comunicação educacional e institucional \\
5 & \texttt{05\_nucleo-andromeda.md} & Redes sociais, engajamento e produção de conteúdo viral \\
6 & \texttt{03\_nucleo-triangulum.md} & Podcasts, YouTube, documentários e jornalismo científico \\
\bottomrule
\end{tabular}

\bigskip


\bigskip
\centerline{\rule{13cm}{0.4pt}}
\bigskip

\subsubsection{Trilha 3 --- Infraestrutura Tecnológica}

\textbf{Para:} Desenvolvedores e equipe técnica.\newline
\textbf{Objetivo:} Entender as ferramentas, sistemas e stack tecnológico do projeto.

\bigskip\noindent
\begin{tabular}{p{\dimexpr 0.333\linewidth-2\tabcolsep}p{\dimexpr 0.333\linewidth-2\tabcolsep}p{\dimexpr 0.333\linewidth-2\tabcolsep}}
\toprule
Ordem & Documento & O que encontrar \\
\hline
1 & \texttt{02\_nucleo-magellan.md} & Portal web, automações, DevOps e infraestrutura digital \\
2 & \texttt{ecossistema.md} & Integração entre os núcleos e fluxo de dados \\
3 & \texttt{02\_Eixos\_Tematicos.md} & Eixo Guaraci --- templates, workflows e automação \\
\bottomrule
\end{tabular}

\bigskip


\bigskip
\centerline{\rule{13cm}{0.4pt}}
\bigskip

\subsubsection{Trilha 4 --- Implementação Operacional}

\textbf{Para:} Equipes executoras do projeto.\newline
\textbf{Objetivo:} Planejar e executar as ações do projeto na prática.

\bigskip\noindent
\begin{tabular}{p{\dimexpr 0.333\linewidth-2\tabcolsep}p{\dimexpr 0.333\linewidth-2\tabcolsep}p{\dimexpr 0.333\linewidth-2\tabcolsep}}
\toprule
Ordem & Documento & O que encontrar \\
\hline
1 & \texttt{fase1.md} & Detalhamento mês a mês da Fase 1 (setup, identidade, lançamento) \\
2 & \texttt{cronograma.md} & Visão geral do cronograma em Gantt \\
3 & \texttt{02\_Eixos\_Tematicos.md} & Organização temática do conteúdo a produzir \\
4 & \texttt{01\_estrutura.md} & Papéis, responsabilidades e reuniões \\
5 & \texttt{parceiros.md} & Editais por instituição com datas e valores de bolsa \\
\bottomrule
\end{tabular}

\bigskip


\bigskip
\centerline{\rule{13cm}{0.4pt}}
\bigskip

\subsubsection{Mapa de Documentos}

\begin{verbatim}
index.md                   \rightarrow Ponto de entrada do projeto
│
├── 01_motivacao.md        \rightarrow Por que o projeto existe
├── 02_contexto -bingo.md   \rightarrow O telescópio e sua missão
├── 01_educomunicacao.md   \rightarrow Fundamentação teórica
├── ecossistema.md         \rightarrow Visão sistêmica dos núcleos
├── 01_estrutura.md        \rightarrow Organograma e governança
│
├── Núcleos Operacionais
│   ├── 02_nucleo -magellan.md      \rightarrow TIC e infraestrutura
│   ├── 04_nucleo -via -lactea.md    \rightarrow Educação e institucional
│   ├── 05_nucleo -andromeda.md     \rightarrow Engajamento digital
│   └── 03_nucleo -triangulum.md   \rightarrow Profundidade audiovisual
│
├── Conteúdo e Planejamento
│   ├── 02_Eixos_Tematicos.md     \rightarrow Temas de conteúdo
│   ├── fase1.md                  \rightarrow Detalhamento da Fase 1
│   └── cronograma.md             \rightarrow Cronograma geral
│
└── Recursos
    ├── parceiros.md              \rightarrow Parceiros e financiadores
    └── orcamento.md              \rightarrow Orçamento consolidado
\end{verbatim}